\begin{table}[!htbp]
\centering
\footnotesize
\setlength{\tabcolsep}{3.5pt}
\caption*{\textbf{Table 4.} Within-district slopes with and without statutory zones, 2026}
\begin{tabularx}{\linewidth}{LCC}
\toprule
Assumptions & Zones ignored: in-zone early ballots may be cast anywhere in the district & Statutory zones: only within the voter's council constituency \\
\midrule
Returns and rolls only & \textbf{[$-$0.378, 0.519]} & \textbf{[$-$0.279, 0.439]} \\
Minimal movement, out-of-zone within 10 points & [0.072, 0.139] & [0.072, 0.138] \\
Uniform movement ($\eta $ = 0.1), out-of-zone within 10 points & \textbf{[$-$0.039, 0.239]} & [0.014, 0.195] \\
Mobility-calibrated movement, out-of-zone within 10 points (D) & \textbf{[$-$0.191, 0.373]} & [0.010, 0.201] \\
Main assumptions (E) & [0.034, 0.174] & [0.054, 0.156] \\
\bottomrule
\end{tabularx}
\par\vspace{2pt}{\raggedright\footnotesize\textit{Notes.} Bold ranges include zero; their sign is not identified. Roll-based share z. The zones are the 155 district-council constituencies, about three neighborhoods each. Rows with uniform movement hold the amount of movement fixed, so the comparison isolates the zone rule; the mobility-calibrated rows recalibrate movement to the zone definition. Adding the council contest's published headcounts changes the statutory-zone ranges by at most 0.003 (Online Appendix Table C2).\par}
\end{table}
